\input format.tex

\usepackage{graphicx}
\graphicspath{{cores/}}

\begin{document}

\vspace*{4mm}
%% 各章节
\setlength{\arrayrulewidth}{1pt}
\fontsize{9.3pt}{11pt}\selectfont
\color{gray2}

\begin{spacing}{1.25}
\begin{LRaside}[.8]{\fontsize{11pt}{11pt}\selectfont 1.2有益菌与有害菌}
\noindent
\includegraphics[width=\linewidth]{youyiyouhaijun.pdf}
\asidebreak %
有益菌，指对人类有益的细菌。它们是肠道里的“正义卫士”，不仅能够产生维生素等有益物质，还能抵御致病菌的入侵，减少炎症。有害菌，是肠道里的“捣蛋鬼”，产生毒素等有害物质。有害菌在肠道内数量较少，一般情况下不会对人体健康产生大的影响，但一旦肠道菌群平衡失调，有害菌异常大量增殖，就可能引发多种健康问题。
\end{LRaside}
\end{spacing}

\vspace*{1mm}
{\noindent\bf 有益菌检测结果\\}
\begin{minipage}[t]{.37\textwidth}
\vspace{-1.5cm}
\hspace{1.5cm}
\includegraphics[scale=0.8]{youyijun_logo.pdf}
\end{minipage}
\hspace{1mm}
\begin{minipage}[t]{.47\textwidth}
\begin{tikzpicture}[scale=1.8,x=1.0cm,y=1.0cm, declare function={
filltop1=11.62;
filltop2=17.61;
fillper=0.7837;}]

\fontsize{8pt}{11pt}\selectfont

\pgfmathparse{{filltop1}}\let\pgfmathresultfilltopa\pgfmathresult
\pgfmathparse{{filltop2}}\let\pgfmathresultfilltopb\pgfmathresult
\pgfmathparse{{fillper}}\let\pgfmathresultfillper\pgfmathresult

\def\filltopa{\pgfmathresultfilltopa}
\def\filltopb{\pgfmathresultfilltopb}
\def\fillpercent{\pgfmathresultfillper}

\pgfmathsetmacro{\filltopaa}{15};

% 直接根据分布求得百分比计算检测值长条覆盖的比例
\pgfmathsetmacro{\filltopbb}{\fillpercent*30};

% 根据均值与检测的比例检测值长条覆盖的比例
%\pgfmathparse{\filltopb/\filltopa}\let\pgfmathswitch\pgfmathresult

%\pgfmathtruncatemacro\one{\pgfmathswitch}
%\pgfmathtruncatemacro\two{2}
%\ifnum\one>\two
%\pgfmathsetmacro{\filltopbb}{30};
%\fi
%\ifnum\one=\two
%\pgfmathsetmacro{\filltopbb}{30};
%\fi
%\ifnum\one<\two
%\pgfmathsetmacro{\filltopbb}{\pgfmathswitch*15};
%\fi

\begin{scope}

\draw[fill=gray1,color=gray1,very thick] (0,0) -- (30mm,0) -- (30mm,2.8mm) -- (0,2.8mm);
\path (0,1.4mm) node [yshift=1.5mm,left=0mm] {平均};
\path (0,1.4mm) node [yshift=-1.5mm,left=0mm] {指数};


\draw[fill=fillcolor1,color=fillcolor1,very thick] (0,0) -- (\filltopaa*1mm,0) -- (\filltopaa*1mm,2.8mm) -- (0,2.8mm);
\path (30mm,1.4mm) node [right=0mm] {\filltopa};
\end{scope}


\begin{scope}[yshift=-5mm]
\draw[fill=gray1,color=gray1,very thick] (0,0) -- (30mm,0) -- (30mm,2.8mm) -- (0,2.8mm);
\path (0,1.4mm) node [yshift=1.5mm,left=0mm] {您的};
\path (0,1.4mm) node [yshift=-1.5mm,left=0mm] {指数};


\draw[fill=fillcolor1,color=fillcolor1,very thick] (0,0) -- (\filltopbb*1mm,0) -- (\filltopbb*1mm,2.8mm) -- (0,2.8mm);
\path (30mm,1.4mm) node [right=0mm] {\filltopb};
\end{scope}

\end{tikzpicture}
\begin{center}
\pdfchatuoliver{xiaoren0}{7}\includegraphics[width=.35cm]{xiaoren08.pdf}
\pdfchatu{xiaoren0}{2}

\textbf{\small 您的有益菌含量高于78.37{\%}的参考人群 }
\end{center}
\end{minipage}


\vspace*{4mm}
\noindent
{\fontsize{9pt}{11pt}\selectfont\bf 主要有益菌结果展示}
\begin{tctabularx}{tabularx={m{4.8cm}<{\centering}m{5.2cm}<{\centering}m{4.55cm}<{\centering}}}
&&
\\[-6pt]
\cellcolor{topcolor} \raisebox{2.614pt}{\color{gray0}\bf 检测项目} &
\cellcolor{topcolor} \raisebox{2.614pt}{\color{gray0}\bf 健康状态提示} &
\cellcolor{topcolor} \raisebox{2.614pt}{\color{gray0}\bf 检测结果评价}
\end{tctabularx}

\vspace*{-4.25mm}
\fontsize{8pt}{11pt}\selectfont
\arrayrulecolor{gray2}
\begin{longtable}{m{4.8cm}m{5.2cm}<{\centering}m{0cm}@{}m{4.61cm}<{\centering}}
\hline
\parbox[c]{\hsize}{\vskip7pt {\lantxh 乳酸杆菌属\\肠道益生菌，能够生成乳酸，抑制有害菌及炎症，调节肠道环境} \vskip7pt} & \parbox[c]{\hsize}{\vskip7pt\centerline{\raisebox{-1.5ex}{\includegraphics[width=5cm,height=0.55cm]{bar.pdf}}}\vskip7pt}  &
\hspace*{-4.83cm}\raisebox{-0.45ex}{\includegraphics{smile.pdf}}
 & \begin{minipage}{4.60cm}\begin{center}{{\lantxh 正常{}} }\end{center} \end{minipage} \\
% & \begin{minipage}{4.60cm}\begin{center}{{\lantxh 正常} }\end{center} \end{minipage} \\
\hline
\parbox[c]{\hsize}{\vskip7pt {\lantxh 双歧杆菌属\\有益菌，降解人体不能消化的多糖，产乳酸，调节免疫及肠道环境} \vskip7pt} & \parbox[c]{\hsize}{\vskip7pt\centerline{\raisebox{-1.5ex}{\includegraphics[width=5cm,height=0.55cm]{bar.pdf}}}\vskip7pt}  &
\hspace*{-4.83cm}\raisebox{-0.45ex}{\includegraphics{smile.pdf}}
 & \begin{minipage}{4.60cm}\begin{center}{{\lantxh 正常{}} }\end{center} \end{minipage} \\
% & \begin{minipage}{4.60cm}\begin{center}{{\lantxh 正常} }\end{center} \end{minipage} \\
\hline
\parbox[c]{\hsize}{\vskip7pt {\lantxh 柔嫩梭菌属\\发酵纤维素产生丁酸等有益物质，抑制肠道炎症，促进肠道健康} \vskip7pt} & \parbox[c]{\hsize}{\vskip7pt\centerline{\raisebox{-1.5ex}{\includegraphics[width=5cm,height=0.55cm]{bar.pdf}}}\vskip7pt}  &
\hspace*{-4.83cm}\raisebox{-0.45ex}{\includegraphics{smile.pdf}}
 & \begin{minipage}{4.60cm}\begin{center}{{\color{green}\lantxh 偏高{\\ \bahao 有利于产生有益物质及抑制肠道炎症}} }\end{center} \end{minipage} \\
% & \begin{minipage}{4.60cm}\begin{center}{{\color{green}\lantxh 偏高 \\ \bahao 有利于产生有益物质及抑制肠道炎症} }\end{center} \end{minipage} \\
\hline
\parbox[c]{\hsize}{\vskip7pt {\lantxh 罗斯拜瑞氏菌属\\产生丁酸等有益物质，抑制肠道炎症，有利于肠道及人体健康} \vskip7pt} & \parbox[c]{\hsize}{\vskip7pt\centerline{\raisebox{-1.5ex}{\includegraphics[width=5cm,height=0.55cm]{bar.pdf}}}\vskip7pt}  &
\hspace*{-3.17cm}\raisebox{-0.45ex}{\includegraphics{cry.pdf}}
 & \begin{minipage}{4.60cm}\begin{center}{{\color{orange}\lantxh 偏低{\\ \bahao 不利于产生有益物质及抑制肠道炎症}} }\end{center} \end{minipage} \\
% & \begin{minipage}{4.60cm}\begin{center}{{\color{red}\lantxh 偏低 \\ \bahao 不利于产生有益物质及抑制肠道炎症} }\end{center} \end{minipage} \\
\hline
\parbox[c]{\hsize}{\vskip7pt {\lantxh 阿克曼氏菌属\\降解粘蛋白、调节免疫，有利于肠黏膜完整性，保持正常体重} \vskip7pt} & \parbox[c]{\hsize}{\vskip7pt\centerline{\raisebox{-1.5ex}{\includegraphics[width=5cm,height=0.55cm]{bar.pdf}}}\vskip7pt}  &
\hspace*{-4.83cm}\raisebox{-0.45ex}{\includegraphics{smile.pdf}}
 & \begin{minipage}{4.60cm}\begin{center}{{\lantxh 正常{}} }\end{center} \end{minipage} \\
% & \begin{minipage}{4.60cm}\begin{center}{{\lantxh 正常} }\end{center} \end{minipage} \\
\hline
\multicolumn{3}{l}{}\\
\multicolumn{3}{l}{\includegraphics[scale=1]{greencircle.pdf} {\fontsize{9.3pt}{11pt}\selectfont\bf 绿色表示健康}   \includegraphics[scale=1]{orangecircle.pdf}  {\fontsize{9.3pt}{11pt}\selectfont\bf 橙色表示需要关注}  \includegraphics[scale=1]{redcircle.pdf} {\fontsize{9.3pt}{11pt}\selectfont\bf 红色表示有风险}}
\end{longtable}

\end{document}
